\documentclass[10pt,a4paper]{amsart}
\usepackage[margin=19mm]{geometry}
\usepackage{amsmath,amssymb,mathtools}
\usepackage[scr=rsfs,cal=euler]{mathalfa}
\usepackage{xfrac}
\usepackage[unicode,hidelinks,bookmarks=false]{hyperref}
\newcommand{\ie}{i.e.\ }
\newcommand{\cf}{cf.\ }
\newcommand{\resp}{respectively\ }
\newcommand{\Subsection}{\subsection}
\providecommand{\nobreakdash}{\nobreak\hbox{-}}
\setlength{\emergencystretch}{2em}
\multlinegap0pt
\usepackage{float}
\usepackage{amssymb}
\usepackage{xcolor}
\newtheorem{theorem}{Theorem}
\title{Circuit Decompositions for Triangulations of Surfaces}
\author{Jens Harlander, Maizie Quatrone}
\date{Source: arXiv:2609.02701v1 (CC BY 4.0)}
\begin{document}
\maketitle

\noindent\emph{Source and licence.} Circuit Decompositions for Triangulations of Surfaces. Version arXiv:2609.02701v1 (CC BY 4.0), distributed by arXiv under the Creative Commons Attribution 4.0 International licence, http://creativecommons.org/licenses/by/4.0/, as recorded in the arXiv licence metadata for this exact version and shown on the abstract page https://arxiv.org/abs/2609.02701v1. Full licence notice: \texttt{licenses/arxiv-2609.02701-CC-BY-4.0.txt}.\par\smallskip

\noindent\textit{Editorial scope.} Counted unit: the article's theorem theorem (source0.tex lines 100--103). The article's complete original proof of it is delivered with it (source0.tex lines 105--138, one \texttt{proof} environment of 34 lines). No mathematics is written, repaired or re-derived: the statement and the proof are the article's own lines, in the article's own notation, and no retained line is substituted or re-flowed.  No further source-local context is needed: every cross-reference of every retained line resolves inside the two delivered spans.  One declared adaptation: exactly 2 ancillary strings inside the retained spans are omitted (image-inclusion commands and, where the source repeats a label, the repeated label definitions) and no other line differs from the source; the figure environments, with their placement, captions and labels, are kept, and the package records the omitted strings in fidelity receipts bound to the affected bodies.  No retained line contains a citation, so the wrapper carries no bibliography and no bibliography entry.  The retained lines contain the article's own diagram code, which the wrapper typesets with the standard tikz packages; no image file is delivered and the package declares no asset dependency.  Editorial formatting only, in addition: an AMS \texttt{amsart} preamble that restates the article's own declarations and macros used by the retained lines, namely the environments \texttt{theorem} on separate counters, together with the standard packages those lines need.  The retained environments print in this document as Theorem 1 in order of appearance, whereas the article prints the same items with the numbering of its own full document.  The complete native source of the article as distributed in the arXiv e-print for this exact version is bundled unchanged as \texttt{sources/209228/source0.tex}.
\begin{theorem}
    If $S$ is an oriented 3-divisible surface, then 
    $$\theta\colon \pi_1(S,v_0)\to \mathbb Z_3,$$ defined by $\theta[p]=\theta(p)\mod3$, is a well defined group homomorphism.
\end{theorem}

\begin{proof} 
    Let $p\in P(S,v_0)$. We have to check that $\theta(p)$ does not change when we apply the homotopy moves 1 and 2 defined at the beginning of this section. Suppose $e_1e_2$ is a subsequence that occurs in $p$. Insert a canceling pair $e_1e\bar ee_2$. We have to check that $\theta(e_1e_2)\equiv\theta(e_1e)+\theta(e\bar e)+\theta(\bar e e_2)\pmod3$. Note that $\theta(e\bar e)=0$, which leaves us with showing $\theta(e_1e_2)\equiv\theta(e_1e)+\theta(\bar e e_2)\pmod3$. There are two cases to consider. Let $v$ be the terminal vertex of $e_1$. First assume that when listing the edges of st$(v)$ that contain $v$ using the orientation at $v$, the edge $e$ appears after $e_1$ but before $e_2$. See Figure \ref{canceling_figure}.
    Then
    $$\theta(e_1e)+\theta(\bar e e_2)=\text{val}(v)+\theta(e_1e_2)\equiv \theta(e_1e_2) \pmod3$$
    \begin{figure}[h!]
        \centering
        
        \caption{This figure shows the insertion of a canceling pair.}
        \label{canceling_figure}
    \end{figure}
    because $S$ is assumed to be 3-divisible, so val$(v)\equiv0\pmod3$. 
    
    \noindent The other case where the edge $e$ appears after $e_2$ is treated in the same way, we leave the details to the reader. This shows that insertion of canceling pairs into the path $p$ does not alter $\theta(p)$.
    
    Nest assume that $e_1e_2e_3$ is a subsequence that occurs in $p$ and we replace $e_2$ by $e_2'e_2''$, where $e_2$, $e_2'$, and $e_2''$ are edges of a single triangle. Let $v$ be the terminal vertex of $e_2'$. We consider the case depicted in Figure \ref{rerouting_figure}. 
    We have 
    \begin{align*}
        \theta(e_1e_2'e_2''e_3) &=\theta(e_1e_2')+\theta(e_2'e_2'')+\theta(e_2''e_3) \\
        &=\theta(e_1e_2)-1+\text{val}(v)-1+\theta(e_2e_3)-1 \\
        &\equiv\theta(e_1e_2)+\theta(e_2e_3)\pmod3
    \end{align*}
    because val$(v)\equiv0\pmod3$. The other case, where the triangle sits on the other side of $e_2$ is treated similarily, and we leave the details to the reader. This shows that rerouting a path across a triangle does not change $\theta(p)$.
    
    \begin{figure}[h]
        \centering
        
        \caption{Rerouting across a triangle. Note that we assume clockwise orientations at all vertices, something we can guarantee only in the orientable case.}
        \label{rerouting_figure}
    \end{figure}
    
    We have shown that $\theta[p]=\theta(p)\mod3$ is well defined. Now notice that 
    $$\theta[pp']\equiv\theta(pp')\equiv\theta(p)+\theta(p')\equiv\theta[p]+\theta[p']\pmod3.$$
    This shows that $\theta$ is a group homomorphism.
\end{proof}

\end{document}
